%% 1-1 いくつ あるかな？（10ページ、p.16〜25）
%% 第1巻のきまり：15 pt・行送り28 pt、分かち書き（文節の間に半角の空白）、漢字はすべてルビ（この章は入学すぐなので、漢字はほとんど使わない）
%% この章で定義した部品（頭文字 bkqa）：\bkqaplate（ふかめの おさら）、\bkqacat（ねこ）、\bkqachair（いす）、
%%   \bkqacartop（上から見た くるま）、\bkqashoe（くつ）、\bkqahand（てのひら）、\bkqaclap（てを たたく）、\bkqabadge（A・B の まる）
\makeatletter
% おさら（\bkplate より ふかく 見える。小さく かいても わかる）  \bkqaplate{(x,y)}{はば}
\newcommand\bkqaplate[2]{\begin{scope}[shift={#1}]
  \fill[bkLine!80] (0,-0.7) ellipse ({#2/2} and {#2/4.2});
  \filldraw[fill=white, draw=bkGray!80, line width=0.8pt] (0,0) ellipse ({#2/2} and {#2/4.2});
  \draw[bkLine, line width=0.6pt] (0,0) ellipse ({#2/2*0.66} and {#2/4.2*0.62});\end{scope}}
% ねこの かお  \bkqacat{(x,y)}{大きさ（1で はば 2mm）}{色}
\newcommand\bkqacat[3]{\begin{scope}[shift={#1}, scale=#2]
  \fill[#3] (-0.95,0.3) -- (-0.78,1.25) -- (-0.2,0.78) -- cycle (0.95,0.3) -- (0.78,1.25) -- (0.2,0.78) -- cycle;
  \fill[#3] (0,0) ellipse (1.0 and 0.85);
  \fill[bkInk] (-0.36,0.1) circle (0.11) (0.36,0.1) circle (0.11);
  \fill[bkAppleRed!70] (0,-0.14) circle (0.08);
  \draw[bkInk!70, line width=0.35pt] (-0.28,-0.2) -- (-1.05,-0.08) (-0.28,-0.3) -- (-1.05,-0.42) (0.28,-0.2) -- (1.05,-0.08) (0.28,-0.3) -- (1.05,-0.42);
\end{scope}}
% いす（よこから）  \bkqachair{(x,y)}{高さ}{色}   (x,y) は あしもとの まん中
\newcommand\bkqachair[3]{\begin{scope}[shift={#1}, scale=#2/10]
  \fill[#3!75!black] (-2.9,0) rectangle (-2.1,4.8) (2.1,0) rectangle (2.9,4.8);
  \fill[#3, rounded corners=0.4] (-3.2,4.6) rectangle (3.2,5.7);
  \fill[#3, rounded corners=0.4] (2.1,5.5) rectangle (3.2,10);
\end{scope}}
% くるま（上から）  \bkqacartop{(x,y)}{色}   (x,y) は まん中。はば 9 mm、ながさ 16 mm
\newcommand\bkqacartop[2]{\begin{scope}[shift={#1}]
  \fill[#2, rounded corners=1.8mm] (-4.5,-8) rectangle (4.5,8);
  \fill[white!75!bkBlue, rounded corners=0.8mm] (-3.4,1.6) rectangle (3.4,4.6) (-3.4,-5.8) rectangle (3.4,-3.8);
  \fill[bkYellow] (-3.7,7.1) rectangle (-2.3,7.8) (2.3,7.1) rectangle (3.7,7.8);
\end{scope}}
% くつ（よこから）  \bkqashoe{(x,y)}{ながさ}{色}   (x,y) は かかとの した
\newcommand\bkqashoe[3]{\begin{scope}[shift={#1}, scale=#2/10]
  \fill[#3, rounded corners=1.2] (0,0) -- (10,0) -- (10,2.3) .. controls (8.4,3) and (6.6,3.3) .. (5.6,4.8) -- (1.2,4.8) -- (0,3.4) -- cycle;
  \fill[white] (0,0) rectangle (10,0.9);
\end{scope}}
% てのひら（ゆび 4本と おやゆび）  \bkqahand{色}   まん中が (0,0)、たかさ 約2
\newcommand\bkqahand[1]{%
  \foreach \fx/\fh in {-0.36/0.8,-0.12/0.95,0.12/0.9,0.36/0.72} {\filldraw[fill=#1, draw=bkStem!70, line width=0.35pt, rounded corners=0.1] (\fx-0.115,0) rectangle (\fx+0.115,0.3+\fh);}
  \filldraw[fill=#1, draw=bkStem!70, line width=0.35pt, rounded corners=0.1, rotate around={35:(-0.45,-0.45)}] (-0.57,-0.45) rectangle (-0.33,0.35);
  \filldraw[fill=#1, draw=bkStem!70, line width=0.35pt, rounded corners=0.3] (-0.5,-0.9) rectangle (0.5,0.38);}
% てを たたく  \bkqaclap{(x,y)}{大きさ（1で たかさ 約3 mm）}
\newcommand\bkqaclap[2]{\begin{scope}[shift={#1}, scale=#2]
  \begin{scope}[shift={(-0.62,-0.35)}, rotate=-20]\bkqahand{bkSkin}\end{scope}
  \begin{scope}[shift={(0.62,-0.35)}, xscale=-1, rotate=-20]\bkqahand{bkSkin!90!bkStem}\end{scope}
  \draw[bkOrange, line width=0.9pt, line cap=round] (-1.5,1.05) -- (-2.05,1.55) (0,1.5) -- (0,2.2) (1.5,1.05) -- (2.05,1.55);
\end{scope}}
% A・B・C の まる  \bkqabadge{(x,y)}{文字}
\newcommand\bkqabadge[2]{\node[circle, fill=bkYellow, text=white, font=\sffamily\bfseries\small, minimum size=5mm, inner sep=0pt] at #1 {#2};}
\makeatother

\bkchapter{1}{1}{いくつ あるかな？}{\bklevel{1}}{いくつ あるかな？}
\bkmichishirube{}{1-2,1-4,1-5}

\begin{goal}
\begin{itemize}[label={\color{bkTheme}●}, leftmargin=1.4\zw, itemsep=0pt, topsep=0pt]
\item 1つずつ ペアに して、 どっちが おおいか わかる。
\item ものを かぞえて、 10までの かずを すうじで かける。
\item リンゴでも くるまでも、 3つなら 「3」と かける。
\end{itemize}
\end{goal}

\begin{hatena}
\begin{bkfigure}
\begin{tikzpicture}[x=1mm, y=1mm]
  \fill[white, rounded corners=3mm] (0,0) rectangle (122,56);
  % テーブル
  \fill[bkStem!25, rounded corners=3mm] (4,4) rectangle (98,52);
  % おさら 5まい（つめて ならぶ）
  \foreach \x in {26,38.5,51,63.5,76} {\bkqaplate{(\x,38)}{11.5}}
  % リンゴ 4こ（はなれて ならぶ）
  \foreach \x in {13,38,63,88} {\bkapple{(\x,16)}{4.6}{bkAppleRed}}
  % こども
  \bkkid{(111,5)}{38}{bkRoute6}{bkHair}
  \node[font=\sffamily\bfseries\fontsize{22}{22}\selectfont, text=bkTheme] at (116,49) {？};
\end{tikzpicture}
\end{bkfigure}
\noindent おやつの じかん。 リンゴを 1こずつ おさらに のせたいな。\\
おさらと リンゴ、 どっちが おおいかな？\\
リンゴの ほうが ながく ならんで いるけど……。
\end{hatena}

\clearpage
%% ---------------- p.17 やってみよう ----------------
\begin{yattemiyou}[せんで むすんで くらべよう]
こどもに ふうせんを 1こずつ わたします。 こどもと ふうせん、 どっちが おおいかな？
\begin{enumerate}[label={\bkcircled[bkGreen]{\arabic*}}, leftmargin=2.2\zw, itemsep=2pt, topsep=2pt]
\item こども 1にんと ふうせん 1こを、 えんぴつで せんで むすぼう。 せんが まがっても だいじょうぶ。
\item あいてが いなくて、 あまった ものは どれかな？
\end{enumerate}
\begin{bkfigure}
\begin{tikzpicture}[x=1mm, y=1mm]
  \fill[white, rounded corners=3mm] (0,-2) rectangle (120,60);
  % こども 4にん（はなれて ならぶ）
  \foreach \x/\c in {15/bkRoute2,45/bkRoute6,75/bkRoute4,105/bkRoute3} {\bkkid{(\x,32)}{25}{\c}{bkHair}
    \fill[bkInk] (\x,28.5) circle (1.1);}
  % ふうせん 5こ（つめて ならぶ）
  \foreach \x/\c in {36/bkAppleRed,48/bkYellow,60/bkBlue,72/bkGreen,84/bkPurple} {\bkballoon{(\x,9)}{4.2}{\c}
    \fill[bkInk] (\x,17) circle (1.1);}
\end{tikzpicture}
\end{bkfigure}
\begin{enumerate}[label={\bkcircled[bkGreen]{\arabic*}}, leftmargin=2.2\zw, itemsep=2pt, topsep=0pt, start=3]
\item おおい ほうの \bkkakikomi[8mm] に ○を かこう。\par\vspace{2pt}
\noindent こども\enspace\bkkakikomi[8mm]\hspace{4\zw}ふうせん\enspace\bkkakikomi[8mm]
\item こんどは かぞえて みよう。\par\vspace{2pt}
\noindent こどもは\enspace\bkkakikomi\enspace にん、\hspace{1.5\zw}ふうせんは\enspace\bkkakikomi\enspace こ
\end{enumerate}
\end{yattemiyou}

\bkfukidashi{mina}{かぞえなくても、 せんで むすぶだけで どっちが おおいか わかったよ！}

\clearpage
%% ---------------- p.18〜19 数学の道具 ----------------
\bksection{doguu}{\R{数学}{すう|がく}の \R{道具}{どう|ぐ}}

\bksubsection{1}{1つずつ ペアに する}

\label{bk:1-1:pea}おさらと リンゴを、 1つずつ せんで むすんで みよう。

\begin{bkfigure}
\begin{tikzpicture}[x=1mm, y=1mm]
  \foreach \x in {26,38.5,51,63.5,76} {\bkqaplate{(\x,25)}{11.5}}
  \foreach \x in {13,38,63,88} {\bkapple{(\x,5)}{4.6}{bkAppleRed}}
  \foreach \a/\p in {13/26,38/38.5,63/51,88/63.5} {\draw[bkBlue, line width=1.3pt, line cap=round] (\a,10.5) -- (\p,21.6);}
  \draw[bkTheme, line width=1.2pt] (76,24.6) ellipse (8.4 and 5.4);
  \draw[bkTheme, line width=1pt] (84.2,26) -- (90,28.5);
  \node[anchor=west, font=\sffamily\bfseries\small, text=bkTheme] at (90,28.5) {あまった！};
\end{tikzpicture}
\end{bkfigure}

\noindent おさらが 1まい あまったね。 \bkhatenaref の こたえは、 「おさらの ほうが おおい」です。\\
1つずつ くみに する ことを、 「ペアに する」と いいます。\bkidx{ぺあ}{ペア}

\noindent ペアに すると、 わかる ことは 3とおりです。

\begin{bkfigure}
\begin{tikzpicture}[x=1mm, y=1mm, lb/.style={font=\sffamily\bfseries\small, text=bkTheme}]
  \foreach \px/\l in {0/A,44/B,88/C} {\fill[white, draw=bkLine, line width=0.8pt, rounded corners=2.5mm] (\px,0) rectangle ++(40,30);
    \bkqabadge{(\px+4.5,25.5)}{\l}}
  % A：3と 3（あまりなし）
  \foreach \x in {11,20,29} {\bkqaplate{(\x,21)}{8}\bkapple{(\x,7)}{3}{bkAppleRed}\draw[bkBlue, line width=1pt] (\x,17.8) -- (\x,10.6);}
  % B：おさら 4、リンゴ 3
  \foreach \x in {51,60,69} {\bkapple{(\x,7)}{3}{bkAppleRed}\draw[bkBlue, line width=1pt] (\x,17.8) -- (\x,10.6);}
  \foreach \x in {51,60,69,78} {\bkqaplate{(\x,21)}{8}}
  \draw[bkTheme, line width=1.1pt] (78,21) circle (5.2);
  % C：おさら 3、リンゴ 4
  \foreach \x in {95,104,113} {\bkqaplate{(\x,21)}{8}\draw[bkBlue, line width=1pt] (\x,17.8) -- (\x,10.6);}
  \foreach \x in {95,104,113,122} {\bkapple{(\x,7)}{3}{bkAppleRed}}
  \draw[bkTheme, line width=1.1pt] (122,7) circle (5);
  \node[lb] at (20,-5) {あまらない};
  \node[lb] at (64,-5) {おさらが あまる};
  \node[lb] at (108,-5) {リンゴが あまる};
  \node[font=\sffamily\bfseries\small] at (20,-12) {→ おなじ かず};
  \node[font=\sffamily\bfseries\small] at (64,-12) {→ おさらが おおい};
  \node[font=\sffamily\bfseries\small] at (108,-12) {→ リンゴが おおい};
\end{tikzpicture}
\end{bkfigure}

\noindent ペアに なった ものどうしは、 おなじ かず。 だから、 あまった ものが ある ほうが、 その ぶんだけ おおいのです。 ならびの ながさは かんけい ありません。

\clearpage
\bksubsection{2}{かぞえて くらべる}

\label{bk:1-1:kazoeru}ゆびで 1つずつ さしながら、 かぞえて みよう。

\begin{bkfigure}
\begin{tikzpicture}[x=1mm, y=1mm,
  num/.style={circle, fill=white, draw=bkBlue, line width=0.8pt, text=bkBlue!80!black, font=\sffamily\bfseries\footnotesize, minimum size=5.6mm, inner sep=0pt},
  last/.style={circle, fill=bkTheme, text=white, font=\sffamily\bfseries\footnotesize, minimum size=5.6mm, inner sep=0pt}]
  \foreach \i in {1,...,4} {\bkapple{(\i*13,23)}{4.4}{bkAppleRed}}
  \foreach \i/\t in {1/いち,2/に,3/さん} {\node[num] at (\i*13,14) {\i};}
  \node[last] at (52,14) {4};
  \node[anchor=west, font=\sffamily\bfseries, text=bkTheme] at (64,18) {→ リンゴは 4こ};
  \foreach \i in {1,...,5} {\bkqaplate{(\i*13,4)}{11.5}}
  \foreach \i in {1,...,4} {\node[num] at (\i*13,-4.5) {\i};}
  \node[last] at (65,-4.5) {5};
  \node[anchor=west, font=\sffamily\bfseries, text=bkTheme] at (77,-0.3) {→ おさらは 5まい};
\end{tikzpicture}
\end{bkfigure}

\noindent かぞえる ときの やくそくは 4つです。
\begin{enumerate}[label={\bkcircled[bkTheme]{\arabic*}}, leftmargin=2.2\zw, itemsep=0pt, topsep=2pt]
\item 1こに、 かずを 1つ いう。
\item とばさない。
\item おなじ ものを 2かい かぞえない。 かぞえた ものに しるしを つけると いいよ。
\item さいごに いった かずが、 ぜんぶの かず。\bkidx{かぞえる}{かぞえる}
\end{enumerate}

\noindent どうして ④に なるのかな？ かぞえる ときは、 「いち、 に、 さん、 し」と いう かずの ことばと、 リンゴを 1つずつ ペアに して います。 ことばが 「し（4）」まで ペアに なったので、 リンゴも 4こ あるのです。

\clearpage
\bksubsection{3}{かずの ならびと \R{大}{おお}きさ}

\label{bk:1-1:narabi}ブロックを 1こずつ ふやして ならべると、 かいだんが できます。

\begin{bkfigure}
\begin{tikzpicture}[x=1mm, y=1mm]
  \foreach \n in {1,...,10} {
    \pgfmathsetmacro\cx{(\n-1)*12.5}
    \ifnum\n=4 \def\cc{bkRoute6}\else\ifnum\n=5 \def\cc{bkAppleRed}\else\def\cc{bkRoute2!70}\fi\fi
    \foreach \k in {1,...,\n} {\filldraw[fill=\cc, draw=white, line width=0.6pt] (\cx,{(\k-1)*4.4}) rectangle ++(9,4.4);}
    \node[font=\sffamily\bfseries\small, text=bkInk] at (\cx+4.5,-4) {\n};}
  \draw[bkTheme, line width=1pt, -{Stealth[length=4pt]}] (42,19.6) .. controls (44,28) and (51,30) .. (54.5,24.4);
  \node[anchor=east, font=\sffamily\bfseries\footnotesize, text=bkTheme] at (40,26) {1つ ふえる};
\end{tikzpicture}
\end{bkfigure}

\noindent となりの かずに すすむ たびに、 1つ ふえます。 4の つぎの 5は、 4より 1つ おおい。 だから、 5は 4より \textbf{おおきい} かずです。 かぞえる とき、 あとから でて くる かずほど おおきいのです。\bkidx{おおきいかず}{おおきい かず}

\noindent \pageref{bk:1-1:kazoeru}ページで、 おさらは 5まい、 リンゴは 4こでした。 5は 4より おおきいので、 かぞえても おさらの ほうが おおいと わかります。 ペアに しても、 かぞえても、 おなじ こたえに なったね。

\bksubsection{4}{どれも 「3」}

\label{bk:1-1:onaji}リンゴ 3こ、 くるま 3だい、 てを たたく 「パン」 3かい。 どれも、 おはじき 3こと ぴったり ペアに なります。

\begin{bkfigure}
\begin{tikzpicture}[x=1mm, y=1mm]
  \foreach \px in {0,44,88} {\fill[white, draw=bkLine, line width=0.8pt, rounded corners=3mm] (\px,0) rectangle ++(40,40);}
  % リンゴ
  \foreach \x in {8,20,32} {\bkapple{(\x,31)}{4}{bkAppleRed}}
  % くるま
  \foreach \x in {52,64,76} {\bkcar{(\x-4.2,27)}{0.36}{bkRoute6}}
  % てを たたく
  \foreach \x in {96,108,120} {\bkqaclap{(\x,29.5)}{2.6}
    \node[font=\sffamily\bfseries\fontsize{9}{9}\selectfont, text=bkOrange!85!black] at (\x,37.2) {パン};}
  % おはじきと せん
  \foreach \px in {0,44,88} {\foreach \d in {8,20,32} {
    \draw[bkGray, line width=0.8pt] (\px+\d,24.5) -- (\px+\d,19.5);
    \bkohajiki{(\px+\d,17)}{2.6}{bkRoute2}}
    \node[font=\sffamily\bfseries\fontsize{24}{24}\selectfont, text=bkTheme] at (\px+20,6.5) {3};}
\end{tikzpicture}
\end{bkfigure}

\noindent くるまは リンゴより ずっと \R{大}{おお}きいね。 てを たたく 「パン」は、 さわる ことも できません。 それでも、 どれも おはじき 3こと ぴったり ペアに なるので、 おなじ かずです。 ものの しゅるいや \R{大}{おお}きさが ちがっても、 おなじ 「3」と いう かずで あらわします。

\bksubsection{5}{すうじを かこう}

\label{bk:1-1:suuji}1から 10までの すうじと よみかた。 10の わくの ● と おなじ かずだよ。\bkidx{すうじ}{すうじ}

\begin{bkfigure}
\begin{tikzpicture}[x=1mm, y=1mm]
  \foreach \n/\r [count=\k from 0] in {1/いち,2/に,3/さん,4/{し・よん},5/ご,6/ろく,7/{しち・なな},8/はち,9/{く・きゅう},10/じゅう} {
    \pgfmathtruncatemacro\cx{mod(\k,5)}\pgfmathtruncatemacro\cy{div(\k,5)}
    \begin{scope}[shift={(\cx*26,-\cy*27)}]
      \fill[white, draw=bkLine, line width=0.8pt, rounded corners=2mm] (0,0) rectangle (24,25);
      \bktenframe{(3.5,16.2)}{3.4}{\n}{bkRoute2}{0}{bkRoute2}
      \node[font=\sffamily\bfseries\fontsize{20}{20}\selectfont, text=bkInk] at (12,10) {\n};
      \node[font=\footnotesize, text=bkInk] at (12,3.4) {\r};
    \end{scope}}
\end{tikzpicture}
\end{bkfigure}

\noindent 10の わくは、 \R{上}{うえ}の だんに 5こ、 \R{下}{した}の だんに 5こ はいります。 5と 5で 10。 10で、 わくが ちょうど いっぱいに なります。\bkidx{じゅうのわく}{10の わく}

\noindent 6から 9は、 \R{上}{うえ}の だんが いっぱいの 5と、 \R{下}{した}の だんの いくつか。 たとえば 7は、 5と 2です。 わくを みると、 1つずつ かぞえなくても かずが わかるよ。 あいて いる ますを かぞえると、 あと いくつで 10に なるかも わかります。

\noindent 10だけは、 「1」と 「0」の 2つの すうじで かきます。 どうして そう かくのかは、 1-2と 1-3で わかるよ。\par\vspace{4pt}

\begin{taisetsu}
1つずつ ペアに して、 あまった ほうが おおい。\\
かぞえる ときは、 さいごに いった かずが ぜんぶの かず。\\
おなじ かずなら、 なんでも おなじ すうじで かく。
\end{taisetsu}

%% （第11回：ここで ページを かえない）
%% ---------------- p.20〜21 例題 ----------------
\begin{reidai}[ペアで くらべる]
1にんに 1こずつ、 あめを くばります。 こどもと あめ、 どっちが おおい？\par\vspace{3pt}
\noindent\begin{minipage}[c]{0.5\linewidth}
\kangaekata むすぶと、 こどもが 1にん あまる。

\kotae こどもが おおい
\end{minipage}\hfill
\begin{minipage}[c]{0.48\linewidth}\centering
\begin{tikzpicture}[x=1mm, y=1mm]
  \foreach \i in {0,...,4} {\bkperson{(4+\i*11,12)}{12}{bkRoute6}}
  \foreach \x in {2,19,36,53} {\bkcandy{(\x,2)}{2.2}{bkRoute4}}
  \foreach \a/\p in {2/4,19/15,36/26,53/37} {\draw[bkBlue, line width=1.1pt, line cap=round] (\a,4.2) -- (\p,11);}
  \draw[bkTheme, line width=1.1pt] (48,18) ellipse (4.2 and 7.6);
\end{tikzpicture}
\end{minipage}
\end{reidai}

\begin{jibun}
ねこと さかな、 どっちが おおい？ せんで むすぼう。

\noindent\begin{minipage}[c]{0.58\linewidth}\centering
\begin{tikzpicture}[x=1mm, y=1mm]
  \foreach \x in {6,24,42} {\bkqacat{(\x,15)}{4.4}{bkOrange!80}\fill[bkInk] (\x,9.6) circle (0.9);}
  \foreach \x in {10,21,32,43} {\bkfish{(\x,0)}{2.8}{bkBlue}\fill[bkInk] (\x,4) circle (0.9);}
\end{tikzpicture}
\end{minipage}\hfill
\begin{minipage}[c]{0.4\linewidth}
{\small おおい ほうに ○}\par\vspace{2pt}
ねこ\enspace\bkkakikomi[8mm]\hspace{1\zw}さかな\enspace\bkkakikomi[8mm]
\end{minipage}
\bkjans{さかなに ○（さかなが 1ぴき あまる）}
\end{jibun}

\begin{reidai}[かぞえて すうじで かく]
はなは なんぼん？\par\vspace{3pt}
\noindent\begin{minipage}[c]{0.5\linewidth}
\kangaekata しるしを つけながら かぞえる。

\kotae 7ほん
\end{minipage}\hfill
\begin{minipage}[c]{0.48\linewidth}\centering
\begin{tikzpicture}[x=1mm, y=1mm]
  \fill[bkLeaf!15, rounded corners=3mm] (-3,-2) rectangle (55,26);
  \foreach \x/\y/\c in {3/12/bkAppleRed,13/1/bkYellow,20/13/bkPurple,29/2/bkAppleRed,36/12/bkRoute6,45/1/bkPurple,48/12/bkYellow} {
    \bkflower{(\x,\y)}{2.4}{\c}
    \draw[bkBlue, line width=1.2pt, line cap=round] (\x+2.4,\y+11.4) -- (\x+3.4,\y+10.4) -- (\x+5.4,\y+13);}
\end{tikzpicture}
\end{minipage}
\end{reidai}

\begin{jibun}
\noindent\begin{minipage}[c]{0.6\linewidth}
おさらの クッキーは なんまい？
\bkkotaeran[6\zw]
\end{minipage}\hfill
\begin{minipage}[c]{0.38\linewidth}\centering
\begin{tikzpicture}[x=1mm, y=1mm]
  \bkqaplate{(22,8)}{45}
  \foreach \x/\y in {7/8,13/3.5,13.5/12,20/7.5,26/12,28.5/3.5,34/9,39/6} {\bkcookie{(\x,\y)}{2.8}}
\end{tikzpicture}
\end{minipage}
\bkjans{8まい}
\end{jibun}

\begin{reidai}[おなじ かずを さがす]
くるま 3だいと おなじ かずは、 A と B の どっち？
\begin{bkfigure}
\begin{tikzpicture}[x=1mm, y=1mm]
  \fill[white, rounded corners=2.5mm] (0,0) rectangle (38,16);
  \foreach \x in {3,14,25} {\bkcar{(\x+1,5)}{0.33}{bkRoute6}}
  \begin{scope}[shift={(44,0)}]
    \fill[white, rounded corners=2.5mm] (0,0) rectangle (36,16);
    \bkqabadge{(4,12)}{A}
    \foreach \x in {11,18,25,32} {\bkballoon{(\x,9)}{2.6}{bkAppleRed}}
  \end{scope}
  \begin{scope}[shift={(84,0)}]
    \fill[white, rounded corners=2.5mm] (0,0) rectangle (36,16);
    \bkqabadge{(4,12)}{B}
    \foreach \x in {14,23,32} {\bkapple{(\x,7.5)}{3.4}{bkAppleGreen}}
  \end{scope}
\end{tikzpicture}
\end{bkfigure}
\kangaekata B の リンゴも 3こ。\kotaenext B
\end{reidai}

\begin{jibun}
てを 4かい たたいたよ。 おなじ かずは どっち？
\begin{bkfigure}
\begin{tikzpicture}[x=1mm, y=1mm]
  \fill[white, rounded corners=2.5mm] (0,0) rectangle (38,16);
  \foreach \x in {6,15,24,33} {\bkqaclap{(\x,6.5)}{2.2}\node[font=\sffamily\bfseries\fontsize{8}{8}\selectfont, text=bkOrange!85!black] at (\x,13.3) {パン};}
  \begin{scope}[shift={(44,0)}]
    \fill[white, rounded corners=2.5mm] (0,0) rectangle (36,16);
    \bkqabadge{(4,12)}{A}
    \foreach \x in {11,18,25,32} {\bkcandy{(\x,6.5)}{1.8}{bkRoute4}}
  \end{scope}
  \begin{scope}[shift={(84,0)}]
    \fill[white, rounded corners=2.5mm] (0,0) rectangle (36,16);
    \bkqabadge{(4,12)}{B}
    \foreach \x in {11,16.5,22,27.5,33} {\bkflower{(\x,2)}{2}{bkPurple}}
  \end{scope}
\end{tikzpicture}
\end{bkfigure}
\bkkotaeran[5\zw]
\bkjans{A（あめ 4こ）}
\end{jibun}

\begin{reidai}[10の わく]
おはじきが 7こ。 あと なんこで いっぱい？\par\vspace{5pt}
\noindent\begin{minipage}[c]{0.56\linewidth}
\kangaekata あいて いる ますは 3つ。

\kotae あと 3こ
\end{minipage}\hfill
\begin{minipage}[c]{0.42\linewidth}\centering
\begin{tikzpicture}[x=1mm, y=1mm]
  \bktenframe{(0,0)}{7}{7}{bkRoute2}{0}{bkRoute2}
  \foreach \x in {17.5,24.5,31.5} {\draw[bkTheme, line width=1pt, dash pattern=on 1.6pt off 1.2pt] (\x,3.5) circle (2.5);}
\end{tikzpicture}
\end{minipage}
\end{reidai}

\begin{jibun}
\noindent\begin{minipage}[c]{0.56\linewidth}
おはじきが 6こ。 あと なんこで いっぱい？
\bkkotaeran[6\zw]
\end{minipage}\hfill
\begin{minipage}[c]{0.42\linewidth}\centering
\begin{tikzpicture}[x=1mm, y=1mm]
  \bktenframe{(0,0)}{7}{6}{bkRoute2}{0}{bkRoute2}
\end{tikzpicture}
\end{minipage}
\bkjans{あと 4こ}
\end{jibun}

%% （第11回：ここで ページを かえない）
%% ---------------- 第11回：えの ない もんだい・かずの 大きさ ----------------
\begin{reidai}[えが ない ときは ○を かく]
いぬが 5ひき、 ほねが 6ぽん。 1ぴきに 1ぽんずつ あげると、 ほねは たりるかな？\par\vspace{3pt}
\noindent\begin{minipage}[c]{0.5\linewidth}
\kangaekata いぬを ○で 5こ、 ほねを ○で 6こ かいて、 せんで むすぶ。 ほねの ○が 1つ あまる。

\kotae たりる（1ぽん あまる）
\end{minipage}\hfill
\begin{minipage}[c]{0.48\linewidth}\centering
\begin{tikzpicture}[x=1mm, y=1mm]
  \node[anchor=east, font=\sffamily\bfseries\footnotesize] at (-2,12) {いぬ};
  \node[anchor=east, font=\sffamily\bfseries\footnotesize] at (-2,0) {ほね};
  \foreach \i in {0,...,4} {\draw[bkRoute6, line width=1pt] (2+\i*8,12) circle (2.6);\draw[bkBlue, line width=1pt] (2+\i*8,9.2) -- (2+\i*8,2.8);}
  \foreach \i in {0,...,5} {\draw[bkStem, line width=1pt] (2+\i*8,0) circle (2.6);}
  \draw[bkTheme, line width=1.1pt] (42,0) circle (4.4);
\end{tikzpicture}
\end{minipage}
\end{reidai}

\begin{jibun}
こどもが 6にん、 いすが 4つ。 1にんが 1つずつ すわると、 みんな すわれるかな？ ○を かいて しらべよう。\par\vspace{2pt}
\noindent\begin{minipage}[c]{0.58\linewidth}
\begin{tikzpicture}[x=1mm, y=1mm]
  \draw[bkLine, line width=0.8pt, rounded corners=2mm, dash pattern=on 2pt off 1.5pt] (0,0) rectangle (66,20);
  \node[anchor=north west, font=\footnotesize, text=bkGray] at (1,19.5) {ここに ○を かこう};
\end{tikzpicture}
\end{minipage}\hfill
\begin{minipage}[c]{0.4\linewidth}
すわれる\enspace\bkkakikomi[8mm]\par\vspace{4pt}
すわれない\enspace\bkkakikomi[8mm]
\end{minipage}
\bkjans{すわれない（2にん すわれない）}
\end{jibun}

\begin{reidai}[どっちが おおきい かず？]
7と 9では、 どっちが おおきい かず？\par\vspace{3pt}
\noindent\begin{minipage}[c]{0.6\linewidth}
\kangaekata かぞえると、 7の あとに 8、 9と でて くる。 あとから でる 9の ほうが おおきい（ブロックも 9の ほうが 2こ おおい）。

\kotae 9
\end{minipage}\hfill
\begin{minipage}[c]{0.38\linewidth}\centering
\begin{tikzpicture}[x=1mm, y=1mm]
  \foreach \k in {1,...,7} {\filldraw[fill=bkRoute2!70, draw=white, line width=0.6pt] (0,{(\k-1)*3.4}) rectangle ++(8,3.4);}
  \foreach \k in {1,...,9} {\filldraw[fill=bkAppleRed, draw=white, line width=0.6pt] (14,{(\k-1)*3.4}) rectangle ++(8,3.4);}
  \node[font=\sffamily\bfseries\small] at (4,-3.6) {7};
  \node[font=\sffamily\bfseries\small] at (18,-3.6) {9};
\end{tikzpicture}
\end{minipage}
\end{reidai}

\begin{jibun}
おおきい ほうの かずに ○を つけよう。\par\vspace{2pt}
\noindent\hspace*{1\zw}（1）\enspace 6\enspace と\enspace 4\hspace{4\zw}（2）\enspace 8\enspace と\enspace 10
\bkjans{（1）6\quad（2）10}
\end{jibun}


\clearpage
%% ---------------- p.22 まちがえやすいところ ----------------
\bksection{tsumazuki}{まちがえやすい ところ}

\begin{tsumazuki}{ながい ほうが おおい？}
\noindent\begin{minipage}[c]{0.5\linewidth}
\bkbatsu\enspace ならびが ながいから、 あかが おおい\par
\bkmaru\enspace むすぶと、 あおが おおい\par
\end{minipage}\hfill
\begin{minipage}[c]{0.46\linewidth}\centering
\begin{tikzpicture}[x=1mm, y=1mm]
  \foreach \x in {0,16,32,48} {\bkohajiki{(\x,11)}{2.6}{bkAppleRed}}
  \foreach \x in {12,19,26,33,40} {\bkohajiki{(\x,0)}{2.6}{bkBlue}}
  \foreach \a/\b in {0/12,16/19,32/26,48/33} {\draw[bkGray, line width=0.8pt] (\a,8.4) -- (\b,2.6);}
\end{tikzpicture}
\end{minipage}
\end{tsumazuki}

\begin{tsumazuki}{すうじの むき}
\bkbatsu\enspace{\sffamily\bfseries\Large\reflectbox{3}\enspace\reflectbox{7}\enspace\reflectbox{9}}\hspace{3\zw}\bkmaru\enspace{\sffamily\bfseries\Large 3\enspace 7\enspace 9}\par
かがみに うつったような むきに ならないように、 \pageref{bk:1-1:suuji}ページの すうじを よく みて かこう。
\end{tsumazuki}

\begin{tsumazuki}{かずを いう こえが おおい}
\noindent\begin{minipage}[c]{0.5\linewidth}
\bkbatsu\enspace 4こなのに 「ご」まで いう\par
\bkmaru\enspace 1こに かずを 1つ\par
\end{minipage}\hfill
\begin{minipage}[c]{0.46\linewidth}\centering
\begin{tikzpicture}[x=1mm, y=1mm]
  \foreach \x in {0,12,24,36} {\bkapple{(\x,0)}{4}{bkAppleRed}}
  \foreach \x/\t in {0/いち,12/に,24/さん,36/し} {\node[font=\sffamily\bfseries\footnotesize, text=bkBlue!80!black] at (\x,8.5) {\t};}
\end{tikzpicture}
\end{minipage}\par\vspace{4pt}
\noindent かずの ことばと ものを、 1つずつ ペアに して かぞえよう。
\end{tsumazuki}

\begin{column}{こいしで かぞえた むかしの ひと}
むかし、 ひつじかいは ひつじ 1ぴきに こいし 1こを ペアに して、 ひつじが みんな かえって きたか たしかめたと いわれて います。
\end{column}

\clearpage
%% ---------------- p.23 つかわれているところ ----------------
\begin{butsuri}
ペアの かんがえは、 まちで はたらいて います。
\begin{bkfigure}
\begin{tikzpicture}[x=1mm, y=1mm]
  % ちゅうしゃじょう（上から）
  \fill[white, rounded corners=3mm] (0,0) rectangle (59,42);
  \fill[bkGray!35] (4,9) rectangle (55,30);
  \foreach \x in {4,14.2,24.4,34.6,44.8,55} {\draw[white, line width=1.2pt] (\x,9) -- (\x,30);}
  \foreach \x/\c in {9.1/bkRoute6,19.3/bkRoute2,39.7/bkRoute4} {\bkqacartop{(\x,18.5)}{\c}}
  \foreach \x in {29.5,49.9} {\fill[bkGreen] (\x,27.6) circle (1.4);}
  \bkdisplay{(15,31.5)}{29}{あき 2}
  \node[font=\sffamily\bfseries\small] at (29.5,4.5) {ちゅうしゃじょう};
  % くつばこ
  \begin{scope}[shift={(63,0)}]
  \fill[white, rounded corners=3mm] (0,0) rectangle (59,42);
  \fill[bkStem!40, rounded corners=1mm] (7,9.5) rectangle (52,39.5);
  \foreach \i in {0,...,3} \foreach \j in {0,1,2} {\fill[bkStem!10] (9+\i*10.8,12+\j*9.6) rectangle ++(9.2,8.2);}
  \foreach \i/\j/\c in {0/0/bkRoute6,1/0/bkRoute2,2/0/bkRoute4,3/0/bkRoute3,0/1/bkRoute2,1/1/bkRoute3,3/1/bkRoute6,0/2/bkRoute4,1/2/bkRoute6,2/2/bkRoute2,3/2/bkRoute3} {
    \bkqashoe{(9.8+\i*10.8,13.4+\j*9.6)}{7.6}{\c}}
  \draw[bkTheme, line width=1.2pt] (9+2*10.8+4.6,12+9.6+4.1) circle (5.6);
  \node[font=\sffamily\bfseries\small] at (29.5,4.5) {くつばこ};
  \end{scope}
\end{tikzpicture}
\end{bkfigure}
\begin{itemize}[label={\color{bkPurple}●}, leftmargin=1.4\zw, itemsep=0pt, topsep=0pt]
\item ちゅうしゃじょうには、 ばしょ ごとに くるまが いるか みる きかいが あります。 ペアに なって いない ばしょを かぞえて、 「あき 2」と みせます。
\item くつばこは 1にんに 1つ。 あいて いる ところで、 まだ きて いない ひとが わかります。
\item きかいは 「1つ きたら 1 ふやす」を、 とても はやく くりかえして かぞえます。 ロボットの モーターも、 まわった かずを かぞえて います。
\end{itemize}
\end{butsuri}

\bkfukidashi{delta}{かぞえる ことは、 たし\R{算}{ざん}の はじまり。 「たしあわせの \R{路線}{ろ|せん}」に のって いこう！}

\begin{bkfigure}
\begin{tikzpicture}[x=1mm, y=1mm]
  \draw[bkRoute2, line width=3pt] (10,0) -- (128,0);
  \node[anchor=west, inner sep=0pt] at (0,12) {\bkrouteicon[1.4em]{2}\hspace{0.3em}{\sffamily\bfseries\small\color{bkRoute2!80!black}たしあわせの \R{路線}{ろ|せん}}};
  \foreach \x/\c/\t/\s in {10/bkVol1/{1-1}/{かぞえる},39.5/bkVol1/{1-4}/{たし\R{算}{ざん}},69/bkVol1/{1-7}/{\R{大}{おお}きな かず},98.5/bkVol1/{1-12}/{かけ\R{算}{ざん}},128/bkVol1/{1-21}/{ひろさ}} {
    \fill[white, draw=\c, line width=1.6pt] (\x,0) circle (2.6);
    \node[anchor=base, font=\sffamily\bfseries\footnotesize, text=\c!85!black] at (\x,-8.4) {\t};
    \node[anchor=base, font=\sffamily\footnotesize, text=bkInk] at (\x,-15.8) {\s};}
  \node[fill=bkTheme, text=white, rounded corners=1mm, font=\sffamily\bfseries\fontsize{8}{9}\selectfont, inner sep=1.5pt] at (10,5.5) {いま ここ};
\end{tikzpicture}
\end{bkfigure}

\clearpage
%% ---------------- p.24〜25 もんだい ----------------
\bkendproblems

\bkproblemgroup{きほん}{1}

\bkproblem いすとりゲームの こどもと いす。 どっちが おおい？
\par\nobreak\vspace{2pt}
\noindent\begin{minipage}[c]{0.66\linewidth}\centering
\begin{tikzpicture}[x=1mm, y=1mm]
  \foreach \x/\c in {0/bkRoute2,14/bkRoute6,28/bkRoute4,42/bkRoute3,56/bkRoute2,70/bkRoute6} {\bkkid{(\x,18)}{13}{\c}{bkHair}\fill[bkInk] (\x,15.8) circle (0.9);}
  \foreach \x in {11,24,37,50,63} {\bkqachair{(\x,0)}{8}{bkStem}\fill[bkInk] (\x,10) circle (0.9);}
\end{tikzpicture}
\end{minipage}\hfill
\begin{minipage}[c]{0.3\linewidth}
{\small せんで むすんで、 おおい ほうに ○}\par\vspace{4pt}
こども\enspace\bkkakikomi[8mm]\par\vspace{4pt}
いす\enspace\bkkakikomi[8mm]
\end{minipage}
\bkpans{こどもに ○（こどもが 1にん あまる）}

\bkproblem なんこ あるかな？ すうじで かこう。
\par\nobreak\vspace{2pt}
\begin{bkfigure}
\begin{tikzpicture}[x=1mm, y=1mm]
  \foreach \px/\l in {0/（1）,43/（2）,86/（3）} {\fill[white, draw=bkLine, line width=0.8pt, rounded corners=2.5mm] (\px,0) rectangle ++(40,22);
    \node[anchor=north west, font=\small, text=bkInk] at (\px+0.5,21.5) {\l};
    \node[inner sep=0pt] at (\px+20,-5) {\bkkakikomi};}
  % （1）リンゴ 6こ
  \foreach \x/\y in {16/16,25.5/17,35/15,10/5.5,20/6.5,30/5} {\bkapple{(\x,\y)}{3.2}{bkAppleRed}}
  % （2）さかな 9ひき
  \foreach \x/\y in {59/17.5,69/18.5,79/16.5,52/11,62/10.5,72/11.5,55/4,66/3.5,77/4.5} {\bkfish{(\x,\y)}{2.2}{bkBlue}}
  % （3）ふうせん 10こ
  \foreach \x/\y in {100/16.5,106/17.5,112/16.5,118/17.5,124/16.5,97/7.5,103/8.5,109/7.5,115/8.5,121/7.5} {\bkballoon{(\x,\y)}{2.2}{bkRoute4}}
\end{tikzpicture}
\end{bkfigure}
\bkpans{（1）6\quad（2）9\quad（3）10}

\bkproblem てを 5かい たたいたよ。 おなじ かずは どれ？
\begin{bkfigure}
\begin{tikzpicture}[x=1mm, y=1mm]
  \foreach \px/\l in {0/A,42/B,84/C} {\fill[white, draw=bkLine, line width=0.8pt, rounded corners=2.5mm] (\px,0) rectangle ++(38,17);
    \bkqabadge{(\px+4.5,12.5)}{\l}}
  % A あめ 4こ
  \foreach \x in {13,20,27,34} {\bkcandy{(\x,6)}{1.8}{bkRoute4}}
  % B リンゴ 6こ
  \foreach \x/\y in {57/11.5,65/11.5,73/11.5,57/4.5,65/4.5,73/4.5} {\bkapple{(\x,\y)}{2.8}{bkAppleRed}}
  % C はな 5ほん
  \foreach \x in {95,101,107,113,119} {\bkflower{(\x,1.5)}{2}{bkPurple}}
\end{tikzpicture}
\end{bkfigure}
\bkkotaeran[5\zw]
\bkpans{C}

\bkproblem 10の わくを いっぱいに するには、 あと なんこ？
\par\nobreak\vspace{2pt}
\begin{bkfigure}
\begin{tikzpicture}[x=1mm, y=1mm]
  \foreach \px/\n/\l in {0/8/（1）,44/5/（2）,88/9/（3）} {
    \node[anchor=west, font=\small, text=bkInk] at (\px-1,5) {\l};
    \bktenframe{(\px+12,0)}{5}{\n}{bkRoute2}{0}{bkRoute2}
    \node[anchor=west, font=\small] at (\px+8,-7.5) {あと\enspace\bkkakikomi\enspace こ};}
\end{tikzpicture}
\end{bkfigure}
\bkpans{（1）あと 2こ\quad（2）あと 5こ\quad（3）あと 1こ}

\bkproblem おおきい ほうの かずに ○を つけよう。\par\vspace{2pt}
\noindent\hspace*{1\zw}（1）\enspace 3\enspace と\enspace 5\hspace{3\zw}（2）\enspace 9\enspace と\enspace 6\hspace{3\zw}（3）\enspace 10\enspace と\enspace 8
\bkpans{（1）5\quad（2）9\quad（3）10}

\bkproblem こうえんの えを みて、 かぞえよう。 かぞえた ものに しるしを つけると いいよ。
\par\nobreak\vspace{2pt}
\begin{bkfigure}
\begin{tikzpicture}[x=1mm, y=1mm]
  \fill[bkLeaf!15, rounded corners=3mm] (0,0) rectangle (128,52);
  \fill[bkBlue!25] (86,10) ellipse (30 and 8);
  \foreach \x/\y in {75/11,86/7,95/12,100/6.5} {\bkfish{(\x,\y)}{2.2}{bkOrange}}
  \foreach \x/\h in {12/22,36/26,118/24} {\bktree{(\x,24)}{\h}}
  \foreach \x/\y/\c in {6/4,18/9,28/3,40/8,52/2,58/14,50/24} {\bkflower{(\x,\y)}{1.9}{bkAppleRed}}
  \foreach \x/\c in {62/bkRoute2,72/bkRoute6,82/bkRoute4,94/bkRoute3,104/bkPurple} {\bkkid{(\x,22)}{17}{\c}{bkHair}}
\end{tikzpicture}
\end{bkfigure}
\noindent\hspace*{1\zw}（1）こども\enspace\bkkakikomi\enspace にん\hspace{1.5\zw}（2）はな\enspace\bkkakikomi\enspace ほん\par\vspace{4pt}
\noindent\hspace*{1\zw}（3）さかな\enspace\bkkakikomi\enspace ひき\hspace{1.5\zw}（4）き\enspace\bkkakikomi\enspace ほん
\bkpans{（1）5にん （2）7ほん （3）4ひき （4）3ぼん}

\bkproblemgroup{チャレンジ}{1}

\bkproblem ちゅうしゃじょうに ばしょが 6つ。 くるまが 7だい きました。 ぜんぶ とめられるかな？ せんで むすんで たしかめよう。
\par\nobreak\vspace{2pt}
\noindent\begin{minipage}[c]{0.62\linewidth}\centering
\begin{tikzpicture}[x=1mm, y=1mm]
  \fill[bkGray!30] (0,0) rectangle (66,15);
  \foreach \i in {0,...,6} {\draw[white, line width=1.1pt] (\i*11,0) -- (\i*11,15);}
  \foreach \i in {0,...,5} {\fill[bkInk] (5.5+\i*11,13) circle (0.9);}
  \foreach \i/\c in {0/bkRoute6,1/bkRoute2,2/bkRoute4,3/bkRoute3,4/bkRoute6,5/bkRoute2,6/bkRoute4} {\bkcar{(\i*9.8-1.5,22)}{0.27}{\c}\fill[bkInk] (\i*9.8+2,20.3) circle (0.9);}
\end{tikzpicture}
\end{minipage}\hfill
\begin{minipage}[c]{0.34\linewidth}
とめられる\enspace\bkkakikomi[8mm]\par\vspace{4pt}
とめられない\enspace\bkkakikomi[8mm]
\end{minipage}
\bkpans{とめられない（1だい あまる）}

\bkproblem こどもが 7にん、 あめが 9こ。 1にんに 1こずつ くばると、 あめは たりるかな？\par
{\small\color{bkOrange!58!black}ヒント：○を 7こと 9こ かいて、 せんで むすんで みよう（れいだい5）。}\par\vspace{2pt}
\noindent\hfill たりる\enspace\bkkakikomi[8mm]\hspace{3\zw}たりない\enspace\bkkakikomi[8mm]
\bkpans{たりる（あめが 2こ あまる）}

\bkproblem 10の わくに おはじきを 7こ いれました。 あと なんこで いっぱいに なるかな？ \R{下}{した}の わくに ○を 7こ かいて しらべよう（れいだい4）。
\par\nobreak\vspace{2pt}
\noindent\begin{minipage}[c]{0.5\linewidth}\centering
\begin{tikzpicture}[x=1mm, y=1mm]
  \bktenframe{(0,0)}{8}{0}{bkRoute2}{0}{bkRoute2}
\end{tikzpicture}
\end{minipage}\hfill
\begin{minipage}[c]{0.46\linewidth}
あと\enspace\bkkakikomi\enspace こ
\end{minipage}
\bkpans{あと 3こ}

\begin{furikaeri}
\begin{checklist}
\item 1つずつ ペアに して くらべられる。\hfill{\footnotesize\sffamily\color{bkTheme}→\pageref{bk:1-1:pea}ページ}
\item かぞえて、 すうじで かける。\hfill{\footnotesize\sffamily\color{bkTheme}→\pageref{bk:1-1:kazoeru}ページ}
\item おなじ かずは、 おなじ すうじ。\hfill{\footnotesize\sffamily\color{bkTheme}→\pageref{bk:1-1:onaji}ページ}
\end{checklist}
チェックが つかなかったら、 →の ページに もどろう。
\end{furikaeri}

\clearpage
%% ---------------- 第11回：ものの かぞえかた（コラム） ----------------
\begin{column}{ものに よって、 かぞえかたが ちがう}
この \R{章}{しょう}では、 「3こ」「3だい」「3まい」と いろいろな いいかたが でて きました。 かずは おなじ 3でも、 ものに よって うしろに つける ことばが かわります。
\begin{bkfigure}
\begin{tikzpicture}[x=1mm, y=1mm, hd/.style={font=\sffamily\bfseries\small, text=bkNavy}, tx/.style={font=\small, anchor=west, text=bkInk}]
  \foreach \k in {0,2,4} {\fill[bkNavy!6] (0,-\k*12-12) rectangle (118,-\k*12);}
  \node[hd] at (14,4) {もの};\node[hd] at (38,4) {つける ことば};\node[hd, anchor=west] at (56,4) {1、 2、 3の いいかた};
  \bkapple{(14,-6)}{3.6}{bkAppleRed}
  \bkpencil{(5,-18)}{18}{bkRoute6}
  \bkqaplate{(14,-30)}{12}
  \bkfish{(14,-42)}{2.6}{bkBlue}
  \bkcar{(8,-58)}{0.4}{bkRoute4}
  \bkkid{(14,-71.5)}{10}{bkRoute2}{bkHair}
  \foreach \w/\r [count=\k from 0] in {こ/{いっこ、 にこ、 さんこ},ほん/{いっぽん、 にほん、 さんぼん},まい/{いちまい、 にまい、 さんまい},ひき/{いっぴき、 にひき、 さんびき},だい/{いちだい、 にだい、 さんだい},にん/{ひとり、 ふたり、 さんにん}} {
    \node[font=\sffamily\bfseries\small, text=bkTheme] at (38,-\k*12-6) {\w};
    \node[tx] at (56,-\k*12-6) {\r};}
\end{tikzpicture}
\end{bkfigure}
\noindent 「ほん」は、 1、 2、 3で 「ぽん・ほん・ぼん」と おとが かわるね。 どの いいかたでも、 かずは かわりません。
\end{column}

\bkfukidashi{sora}{おうちの ひとと、 おさらと おはしを 1つずつ ペアに して ならべて みよう！}

\clearpage
